\section{The rank and the marked lattice on a column}

% BEGIN H77S1
\paragraph{H77S1. The stable range is built into the context.}
Let $K$ be an original context with scale $s$. Let $X$ be any of the
stages over $K$, in the precise sense $K.\mathrm{Over}(X.B.\mathrm{toPre})$.
Write $D=X.\mathrm{toD7}$, and abbreviate its rank parameters by
$f,r,t,n_V$. Unpacking the witness of being over $K$ identifies
the ranks of $X.B$ with the ranks in $K.\mathrm{ok.rank}$. The last
equality of this rank specification is
\[
 n_V=r+t=8fr+4+\mathrm{TOv}(s)\,\mathrm{nWv}(D.R,s).
\]
Both factors of the last product are natural numbers. This proves
$8fr+4\le n_V$ immediately, for every allowed choice of the later
stages. In particular the restriction contributes no new scale
threshold to Q76Z1.
% END H77S1

% BEGIN H77S2
\paragraph{H77S2. The coordinate lattice at the definite space.}
For $0\le i<m$ let $V$ be the source's $V_c(D,i)$, the space
at index $i+1$, with its specified ordered frame $(e_a)$.
Being over $K$ identifies the spaces with the diagonal construction.
The sign at $i+1$ is positive by the base sign condition.
In the defining model the Gram entries are $\delta_{ab}$ and
$L_0$ is the $\mathcal O_E$-span of these very basis vectors.
Transporting the description along the specified equality of spaces
preserves both statements. Uniqueness of coordinates therefore gives
\[
  \sum_a x_a e_a\in L_0
     \quad\Longleftrightarrow\quad x_a\in\mathcal O_E\quad\text{for all }a.
\]
In the definition of $\mathrm{VStd}$ the lattice at the identity
adele is read back from its finite completions. Here a rational
coordinate belongs to every completed ring of integers exactly
when its prime ideal valuations are all nonnegative. Indeed its
principal fractional ideal then lies in $\mathcal O_E$. Applying
this to each coordinate identifies the recovered lattice with
$\mathcal O_E^{\operatorname{rank} V}$. Together with the Gram
calculation this gives the native $\mathrm{VStd}(D,i)$, simultaneously
for every row of $X$. Merely prescribing the Gram entries would omit
the required identification of $L_0$.
% END H77S2

% BEGIN H77S3
\paragraph{H77S3. Selected finite places and the base half.}
Let $w$ be a pinned column. For each selected prime $v$ of the real
field $F$, its chosen prime $P=\mathrm{sel}(v)$ of the CM field $E$
contracts to $v$. If $P$ and $P'$ occur in the tuple prime sets and
lie over the same rational prime $q$, the descriptions of those sets
give $v,v'$ with $P=\mathrm{sel}(v)$ and $P'=\mathrm{sel}(v')$.
Contraction gives $q\in v$ and $q\in v'$. The column's half pin at
$v,v',q$ identifies the primes below $P,P'$ in the host field
$K_{\rm host}$. Taking their ideals gives
$P\cap\mathcal O_{K_{\rm host}}=P'\cap\mathcal O_{K_{\rm host}}$,
the equality required by $\mathrm{SelHalf}$.
Also put $W_t=\{v:\operatorname{char} k(v)\in T.\mathrm{pl}\}$
on the column's place type. Above each member of the finite set
$T.\mathrm{pl}$ there are only finitely many primes, by factorization
in the ring of integers. The conjugate of a prime has the same
residue characteristic. Thus this particular $W_t$ is finite and
conjugation-stable. The used tuple type is finite too, being indexed
by the base tuple together with the finite set of used tuples.
These verifications take place before any analytic record is chosen.
% END H77S3
